\section{Methodology}

\subsection{Cohort Construction and Prediction Task}
\label{sec:cohort}

We build a cohort of 33{,}655 patients from MIMIC-IV, a de-identified electronic health record database from Beth Israel Deaconess Medical Center (BIDMC), an academic medical center in Boston. Each patient's \emph{index date} is set to their earliest hospital admission associated with a cardiometabolic condition: type 2 diabetes, hypertension, hyperlipidemia, obesity, or metabolic syndrome. The index date is the reference point for everything that follows: it marks the boundary between information the model is allowed to see and the outcome it is asked to predict.

Two conditions must both hold for a patient to be included. First, the patient must have at least two hospital visits in total, so that the cohort is not dominated by single-encounter patients whose subsequent course is unknown. Second, among those patients, we require either 90 days of follow-up after the index date or an observed target outcome, whichever comes first; a patient who is only followed for a few days and never develops one of the target diseases is dropped, because there is no way to tell whether they were truly disease-free or simply lost to observation too soon. A patient who develops a target disease before 90 days is kept regardless, since the outcome itself is what we are trying to predict, and it is already known at that point.

We apply a washout step on top of this: patients with a documented history of any of the five target circulatory diseases on or before the index date are excluded. Without this step, a patient's disease could be recorded again during a later admission and mistakenly counted as a new occurrence, when it is really a pre-existing condition already present at the start of observation.

Starting from the index date, we predict the first occurrence of one of five circulatory outcomes: myocardial infarction (MI), ischemic stroke, heart failure (HF), atrial fibrillation (AF), or peripheral artery disease (PAD), and whether that occurrence falls within 1, 3, or 5 years. A patient can, in principle, go on to develop more than one of these five diseases over time, but only the first one that occurs is used as their outcome; developing one of these diseases first also rules out observing another as the patient's first event, so we treat this as a competing-risks survival problem rather than five independent binary classification tasks. Patients who develop none of the five diseases are followed until their last available clinical record and are administratively right-censored at that point, meaning their true disease status after that point is simply unknown, not assumed to be negative.

An occurrence only counts as an outcome event when the disease is recorded as the \emph{primary diagnosis} for a hospital admission, rather than appearing anywhere in the chart. This is meant to separate a genuinely new disease episode, serious enough to be the reason for that admission, from an incidental or historical mention of the same diagnosis recorded during an unrelated visit.

It is worth being precise about five terms that recur throughout this paper, since they are easy to conflate. The \emph{index date} is the single reference day described above. The \emph{pre-index observation window} is the period before the index date from which a model is allowed to draw input information, defined precisely in the Temporal Clinical Representation subsection below. The \emph{prediction horizon} is the period after the index date over which an outcome is assessed (1, 3, or 5 years). An \emph{event} is the first target disease diagnosed as a primary diagnosis after the index date; a \emph{competing event} is any of the other four target diseases occurring first instead of the one currently being evaluated; and \emph{administrative censoring} is the case where a patient's last available record ends without any target disease being observed, leaving their eventual status unknown rather than negative. No model in this study has access to any information dated at or after a patient's own index date.

\subsection{Temporal Clinical Representation}
\label{sec:tkg}

Each patient's clinical history is represented as a collection of timestamped medical facts drawn from structured EHR data. We use observations from the five years before the index date, across eight data categories: diagnoses, procedures, medication orders, laboratory results, ICU stays, vital signs, outpatient measurements, and intravenous fluid administration. Each observation is written as

\[
f_i = (p_i, c_i, t_i, v_i),
\]

where $p_i$ identifies the patient, $c_i$ identifies the medical concept observed (an ICD-9/ICD-10 code for diagnoses and procedures, a normalized drug name for medications, or a standardized identifier for labs and vitals), $t_i$ is the time the observation was recorded, and $v_i$ is an optional numeric value attached to the observation. The value field is populated for observations such as lab results and vital signs, and left undefined for observations that do not carry a numeric value, such as a diagnosis code on its own.

Across the cohort, this produces approximately 14.5 million timestamped facts spanning approximately 33{,}000 distinct medical concepts. This is related to, but not the same as, a conventional temporal knowledge graph (TKG). A standard TKG represents a fact as a quadruple $(s, r, o, t)$: two entities $s$ and $o$ connected by a relation $r$ at time $t$. What we build at this stage is a set of patient-concept observations only; no relation between two concepts, and no graph structure of any kind, is imposed here. This representation is a fact store, not yet a graph. Graph structure, where it exists in this study, is added later by one specific model, described below, not by this shared representation.

To prevent information leakage, only observations from the five years before a patient's own index date are included, and this restriction is enforced once, at the point where this representation is built, before any model-specific processing happens. Because the cutoff is applied here rather than separately inside each model, no observation dated at or after a patient's index date can reach any of the four models described below, regardless of how each model subsequently transforms the data.

This shared representation is the common starting point for all four models, but each one transforms it differently. Cox regression and XGBoost collapse it into a fixed-length feature vector per patient. The TKG-Transformer keeps the events in temporal order and processes them as a sequence. The patient-graph GNN turns the same facts into an explicit graph with patient and medical-concept nodes. These four transformations are described next.

\subsection{Predictive Models}
\label{sec:models}

\textbf{Cox Regression.}
For each patient, we build a fixed-length feature vector from their pre-index history, with approximately 17{,}700 features in total. About 17{,}400 of these are binary bag-of-concepts features: one indicator per medical concept observed anywhere in the training set, marking whether that concept occurs at least once in this patient's pre-index record. Concepts that appear only in the validation or test data, and never in training, are not part of this feature space; an observation of such a concept simply contributes nothing, since no column exists for it. For concepts tied to a numeric measurement, mainly labs and vitals, we add six summary statistics: mean, minimum, maximum, most recent value, observation count, and the slope of value against time. These summary features add 246 dimensions and describe both the typical level and the trend of a patient's numeric measurements over time. The remaining five features are static patient-level covariates: age, sex, Charlson Comorbidity Index, number of baseline cardiometabolic conditions, and whether the patient had a pre-index ICU stay.

We use elastic-net regularized Cox proportional-hazards regression as a linear survival baseline. For covariates $\mathbf{x}$, the cause-specific hazard for disease $c$ is modeled as

\[
h_c(t \mid \mathbf{x}) = h_{0c}(t)\exp(\boldsymbol{\beta}_c^\top\mathbf{x}),
\]

where $h_{0c}(t)$ is the baseline hazard for disease $c$ and $\boldsymbol{\beta}_c$ is the vector of learned coefficients for that disease. We fit five separate Cox models, one per target disease. When fitting the model for disease $c$, a patient whose first event is one of the other four target diseases has their follow-up for disease $c$ end at that point and is treated as a competing event, not as a patient who survived disease $c$.

With roughly 17{,}700 features and comparatively few observed outcomes, this feature space is far larger than the number of events available to fit it, so we regularize. We use elastic-net regularization with an $\ell_1$ ratio of $0.9$, which drives most coefficients to exactly zero while keeping a smaller $\ell_2$ component that lets the model retain groups of correlated features together rather than arbitrarily picking one and discarding the rest. The strength of this regularization is chosen separately for each disease using the training and validation sets, following the procedure described in the Model Training and Hyperparameter Tuning subsection below.

\textbf{XGBoost.}
XGBoost serves as a nonlinear survival baseline and receives exactly the same fixed-length patient representation used for Cox regression. Giving both models identical input means that any difference in their behavior reflects a difference between the models themselves, not a difference in what information each one had access to. As with Cox, we fit a separate model per target disease.

Cox regression assumes that each feature has a fixed, linear effect on the log-hazard. XGBoost instead represents the hazard using an ensemble of gradient-boosted decision trees, which can represent nonlinear effects and interactions between features that a linear model cannot. We train each disease-specific model using XGBoost's Cox partial-likelihood objective (\texttt{survival:cox}). The survival target for each model is encoded as $+t$ when the target disease is observed at duration $t$, and $-t$ when the patient's follow-up ends without that disease occurring, whether because their available record simply ends or because one of the other four target diseases occurs first. The resulting model outputs a relative risk score, where a higher score corresponds to greater estimated risk of the target disease; this is an architectural difference from Cox, and on its own does not indicate which model will perform better.

The number of trees, tree depth, learning rate, subsampling ratio, and feature-subsampling ratio are tuned as described in the Model Training and Hyperparameter Tuning subsection below.

\textbf{TKG-Transformer.}
Despite its name, the TKG-Transformer does not perform any message passing between patients and does not operate on a graph. The name reflects that its input is drawn from the same temporal fact store described above; the model itself processes one patient's own sequence of events in isolation, with no connectivity to any other patient. This is the key architectural difference from the patient-graph GNN described next.

For each patient, we keep the 256 most recent pre-index events, truncating longer histories and padding shorter ones. Each event in the sequence is represented by three combined components: a learned embedding of the medical concept, a continuous-time embedding of how long before the index date the event occurred, and an embedding of its numeric value together with an indicator of whether a value applies to that event (an event such as a diagnosis code, which carries no numeric value, is marked accordingly rather than assigned an arbitrary placeholder number). Time is represented using a Bochner/TGAT-style sinusoidal encoding of the exact, continuous time gap to the index date, rather than sorting events into discrete time bins, so the model can distinguish events that occurred close together in time.

The resulting sequence is processed by a two-layer Transformer encoder with four self-attention heads and a 128-dimensional hidden state. Self-attention lets the model weigh relationships between events at different points in a patient's history while keeping their temporal position. The resulting patient representation is passed to a DeepHit-style competing-risks prediction head, which outputs a probability over every combination of target disease and a discrete time bin; these probabilities are accumulated across time bins to produce a cause-specific risk estimate at the 1-, 3-, and 5-year horizons for each of the five diseases.

Both the TKG-Transformer and the patient-graph GNN train with early stopping, but neither model is allowed to select a checkpoint as its final, "best" model before 15 full passes through the training data. This floor was added after we observed, across all five training seeds, that early stopping consistently selected a checkpoint from epoch 1 or 2 by validation AUROC, and a fidelity check on that checkpoint, described under Explanation Fidelity below, found that the events it flagged as important were statistically indistinguishable from a randomly chosen set, even though its raw discrimination was already competitive. The 15-epoch floor makes any checkpoint before that point ineligible for selection, regardless of its validation score.

\textbf{Patient-Graph GNN.}
The patient-graph GNN is the one model in this study that builds an explicit graph and allows information to move between different patients. It transforms the same pre-index facts described above into a heterogeneous graph with two node types: 33{,}655 patient nodes and 35{,}501 medical-concept nodes, for 69{,}156 nodes in total.

A patient is connected to a concept if that concept was observed in their pre-index history. Each patient-concept pair collapses into a single edge: if a patient has multiple observations of the same concept, only the most recent one determines that edge, which is then assigned to one of three relation types by how recently it occurred relative to the index date: within 90 days, between 90 and 730 days, or more than 730 days before. This means a concept observed only in the distant past, if it was also observed more recently, is represented by one recent-window edge rather than by a separate edge in each window it happened to fall into. Every patient-concept edge is added in both directions (patient-to-concept and concept-to-patient, as separate relation types), so information can flow either way along it.

We add two further sources of concept-to-concept edges. Ontology edges connect a concept to related concepts in a manually built medical hierarchy, such as a diagnosis code to its diagnostic category or a medication to its drug class; these are added in both directions as well. Co-occurrence edges connect concepts that frequently appear together in the same patient's history; these are symmetric and are built using only training-set patients, so no validation or test information influences which concepts are linked. Altogether, the graph has approximately 2.7 million directed edges across 19 relation types.

We process this graph with a relational graph convolutional network (R-GCN) using basis decomposition: a 128-dimensional hidden state, four relation bases, and two message-passing layers, for approximately 4.7 million trainable parameters. With two layers of message passing, a patient node's final representation can incorporate information from its own directly connected concepts (one step of propagation), and, through those same concepts, information associated with other patients connected to the same concepts (a second step of propagation). This is what allows the patient-graph GNN to draw on patients with related clinical histories, which the purely patient-specific TKG-Transformer has no mechanism to do. The resulting patient embeddings are passed to the same DeepHit-style competing-risks head used for the TKG-Transformer. Training is full-batch: one forward and backward pass over the entire graph per epoch, rather than training on batches of individual patients.

Because every cohort patient, including those in the validation and test sets, is represented as a node in this one shared graph, the primary model is trained \emph{transductively}. Concretely, this means: validation and test patients are present in the graph during training, and their allowed pre-index facts and connectivity can influence how information propagates through the graph; but their outcome labels are never used to compute the training loss or update any model parameter, and no information dated at or after a patient's own index date is present anywhere in the graph for any patient. To check whether this transductive setup itself affects performance, we also train and evaluate a second, inductive version of the model in which all validation and test patient edges are removed during training, so those patients contribute no connectivity at all until evaluation.

\subsection{Model Training and Hyperparameter Tuning}
\label{sec:tuning}

Every hyperparameter and every checkpoint is selected using only the training and validation sets. The test set is not used for any selection decision and is evaluated exactly once, after all modeling choices are finalized.

For Cox regression, we fit the entire elastic-net regularization path in a single call, evaluating 100 candidate $\alpha$ values from strongest to weakest penalty, separately for each disease. For each disease, we score every candidate $\alpha$ on the validation set and keep the one with the highest mean AUROC averaged across the 1-, 3-, and 5-year horizons; the final model at test time uses this per-disease selected $\alpha$.

For XGBoost, we tune the number of trees, learning rate, maximum tree depth, subsampling ratio, and feature-subsampling ratio using a coordinate-wise search: starting from a fixed default configuration, we vary one hyperparameter at a time over a small candidate list, keep whichever value gives the best score, and combine the best value from each axis into a single final configuration. Unlike Cox's per-disease $\alpha$, this configuration is shared across all five disease-specific XGBoost models; a single set of hyperparameters is selected once, using the mean validation AUROC at the 3-year horizon averaged across the five diseases.

The TKG-Transformer and the patient-graph GNN train with early stopping on validation performance, subject to the 15-epoch minimum described in the Predictive Models subsection above. For the patient-graph GNN specifically, we additionally tune the hidden dimension, number of message-passing layers, number of relation bases, dropout rate, and learning rate together, using a bounded random search over the joint hyperparameter space (12 randomly sampled configurations plus the existing default, drawn from hidden dimension $\{64,128,256\}$, layers $\{1,2,3\}$, relation bases $\{2,4,8,16\}$, dropout $\{0.0,0.15,0.3\}$, and learning rate $\{3\times10^{-4}, 10^{-3}, 3\times10^{-3}\}$), rather than tuning each hyperparameter independently. We used a joint search because an earlier attempt at tuning each hyperparameter on its own found that combining the single best value from each axis produced a less stable model than the untuned defaults, which pointed to interactions between these hyperparameters that independent, one-at-a-time tuning cannot capture. Configurations were screened by validation mean AUROC at the 3-year horizon, and the best configuration was then confirmed across all five training seeds; this search selected the same configuration already in use as the model's defaults. Every candidate configuration in every search above was evaluated on validation data only.

\subsection{Evaluation}
\label{sec:eval}

\subsubsection{Horizon-Specific Discrimination}
\label{sec:eval-auroc}

We measure discrimination using the area under the receiver operating characteristic curve (AUROC), computed separately for each of the five diseases at each of the three prediction horizons $h \in \{1, 3, 5\}$ years. For disease $c$ at horizon $h$, a patient is labeled positive if $c$ is their first observed event and it occurred within $h$ years of the index date. A patient is labeled negative if they either remained event-free through $h$ years, or developed a different target disease before $h$ years; in the second case, that patient still counts as a negative for disease $c$ at this horizon, rather than being removed from the calculation, since a competing event resolves their status for disease $c$: they are confirmed not to have had it within that window. A patient who was administratively censored before $h$ years, with no target disease observed by their last record, is excluded from the evaluation at that horizon rather than labeled either way, since their status is genuinely unknown. This labeling is independent of, and defined separately from, the cause-specific censoring rule used when fitting Cox: at evaluation time, only true competing events count as negatives, while unresolved censoring is excluded rather than counted.

\subsubsection{Multi-Seed Evaluation}

XGBoost, the TKG-Transformer, and the patient-graph GNN are each trained five times, from five different random seeds, to account for the randomness in their training procedures (parameter initialization and, depending on the model, data ordering or subsampling). For these three models, we report the mean and standard deviation of AUROC across the five resulting trained instances. Cox regression is fit only once: given a fixed design matrix and a fixed regularization strength, its solution is deterministic, so repeating it would not add information.

\subsubsection{Statistical Comparisons}
\label{sec:stats}

We compare models using two different statistical tests, applied depending on what is being compared, because neither one alone answers the full question.

When a stochastic model (trained across five seeds) is compared against Cox, which is deterministic, we use a one-sample $t$-test: the five seed-level AUROC values are treated as a sample and tested against Cox's single, fixed AUROC value. This test asks whether Cox's value falls outside the range of variation the stochastic model produces on its own across training runs. When two stochastic models are compared against each other (for example, the TKG-Transformer against XGBoost), we instead use a two-sample test comparing their two sets of five seed-level AUROC values directly. In both cases, comparisons are Bonferroni-corrected within their own family (the set of comparisons sharing the same pair of models being compared).

Neither of these two tests checks whether the two models actually disagree on the same patients, and neither accounts for the comparator's own sampling uncertainty on this specific test set. For that, we separately apply DeLong's test, a paired test for correlated AUROC estimates computed from two models' predictions on the same test patients. We run DeLong's test once per seed: for a stochastic model with five trained instances, this gives five separate paired comparisons against the same fixed comparator, each on the same test set.

We call a comparison against Cox \emph{robust} only when it satisfies all three of the following: (1) the Bonferroni-corrected one-sample $t$-test on the seed-level AUROC values is significant; (2) the paired DeLong test is significant in at least four of the five seeds; and (3) all five seeds agree on the direction of the difference. A comparison that satisfies only one of these criteria is reported as such, not as robust. As a sensitivity check, we additionally compare two ways of correcting for the combined set of comparisons made by the TKG-Transformer and the patient-graph GNN together (60 comparisons in total, 30 from each model against the two baselines): a single Bonferroni correction applied across all 60 at once, and a Benjamini-Hochberg false-discovery-rate correction applied across the same 60. Confidence intervals for individual AUROC and AUPRC values are computed using a 2{,}000-resample patient-level bootstrap with the percentile method.

\subsubsection{Performance by Pre-Index History Density}
\label{sec:thinrich}

We separately test whether the patient-graph GNN's relative performance against Cox depends on how much pre-index clinical history a patient has, since this is a condition under which patient-to-patient message passing is, in principle, positioned to help. Using only training-set patients, we compute the median number of pre-index facts per patient, and use this training-set median as a fixed cutoff applied to all splits. Test patients below this cutoff are labeled "thin," and those at or above it are labeled "rich." Within each of these two groups, and separately for each of the five trained seeds, we compute AUROC for the patient-graph GNN and for Cox, and take their difference. For a given disease and horizon, we count the pattern as confirmed if the thin-group margin exceeds the rich-group margin in at least four of the five seeds, and we summarize the result across all 15 disease-horizon cells with a sign test. This analysis tests whether relative performance varies with history density; it does not, on its own, establish that message passing is the mechanism responsible for any pattern observed.

\subsubsection{F1 and AUPRC}
\label{sec:f1auprc}

AUROC and AUPRC are both threshold-free, but F1 requires choosing a decision threshold, and with target-event prevalence this low, that choice matters. We report the best achievable F1: for a given set of predictions, F1 is computed at every threshold along the precision-recall curve, and the maximum value obtained is reported, rather than picking one threshold in advance and applying it universally. We do this because recall alone is degenerate at this prevalence: labeling every patient positive trivially achieves a recall of 1.0, so recall by itself cannot distinguish between hyperparameter configurations, and in practice every configuration we tried for XGBoost and the patient-graph GNN tied at exactly that value.

Because of this, we ran a separate hyperparameter search for XGBoost and the patient-graph GNN that selects configurations by validation best-achievable F1 at the 3-year horizon, rather than by AUROC, following the same tuning discipline described above under Model Training and Hyperparameter Tuning (train/validation only, test touched once at the end). Cox has no equivalent F1-specific search, since its only tunable hyperparameter is the elastic-net regularization strength already selected above; its existing, AUROC-tuned risk scores are used directly for this comparison.

Using the resulting best-F1 configurations, we compare the patient-graph GNN against each baseline with two procedures. The first is a paired bootstrap test (2{,}000 resamples): test patients are resampled, and each model's best-F1 threshold is recomputed freshly within every resample, rather than fixed at its original value, so that the uncertainty from threshold selection itself is included in the resulting confidence interval, not assumed away. The second is a one-sample $t$-test comparing the patient-graph GNN's five seed-level best-F1 values against the baseline's single best-F1 value, following the same one-sample-versus-fixed-comparator logic described under Statistical Comparisons above. AUPRC is computed directly, without a threshold, and its confidence intervals are computed with the same 2{,}000-resample bootstrap procedure used for AUROC.

\subsubsection{Explanation Fidelity}
\label{sec:fidelity}

For the TKG-Transformer and the patient-graph GNN, we test whether the individual facts each model's explanation method flags as important actually matter to that model's own prediction, rather than only appearing plausible. We do this with two measures, comprehensiveness and sufficiency, both defined in terms of how much a model's prediction changes when specific input facts are removed or kept.

Comprehensiveness asks: if we remove the facts flagged as important, does the model's prediction change more than if we had removed an equally sized random set of facts instead? A larger change indicates the flagged facts were doing more work. Sufficiency asks the reverse question: if we keep only the flagged facts and remove everything else, does the model's prediction stay closer to its original value than if we had kept only a random set of the same size? A smaller change indicates the flagged facts were, on their own, enough to reproduce roughly the same prediction. In both cases, we quantify "how much the prediction changes" using the KL divergence between the model's original output distribution over disease and time bin, and its output distribution after the facts in question are removed or restricted. For each patient, the same procedure is applied to a random set of facts of the same size, giving a matched, patient-specific baseline for comparison. The number of facts selected, top or random, is 20\% of that patient's own total pre-index event or concept count (with a minimum of one).

The two models identify important facts differently, and it is worth being precise about the difference, since one is not simply an application of the other. For the TKG-Transformer, we use PyTorch Geometric's implementation of GNNExplainer directly: each patient's event sequence is represented to the explainer as a graph with one node per event and no edges between them, and GNNExplainer optimizes a single necessity score per node (one scalar mask value per event) over 100 epochs to best explain the model's full output distribution across all diseases and time bins. For the patient-graph GNN, GNNExplainer's usual per-patient, iterative mask-optimization procedure is computationally impractical, because this model makes one shared forward pass over the entire graph for every patient at once rather than a separate forward pass per patient. We instead use a single gradient-based saliency score per patient: the gradient of that patient's predicted cumulative-incidence probability, for their own actual disease outcome at the 3-year horizon, with respect to the shared concept-embedding table, restricted to the concepts that patient is connected to. This gradient magnitude ranks that patient's connected concepts by importance. Although the ranking signal used for the patient-graph GNN targets a single disease and horizon rather than the full output distribution used for the TKG-Transformer, the fidelity evaluation itself is computed identically for both models, using each model's full disease-and-time-bin output distribution before and after the perturbation.

For each of up to 30 true-positive test patients per disease (150 patients in total across the five diseases), we compare the top-ranked facts against a same-size random set on both comprehensiveness and sufficiency, using a paired $t$-test, a Wilcoxon signed-rank test, and a win rate: the fraction of patients for whom the flagged facts land on the expected side of the random set for that metric (a larger change for comprehensiveness, a smaller change for sufficiency), with a bootstrap 95\% confidence interval on that fraction computed from 5{,}000 resamples. We report both the paired $t$-test and the Wilcoxon test together, rather than only one, because they can disagree: the $t$-test is sensitive to the mean effect across patients and can be driven by a subset of patients with a large effect, while the Wilcoxon test is more sensitive to whether the effect holds for a typical patient. When the two tests disagree, that disagreement is itself part of the result.

\subsubsection{Calibration}
\label{sec:calibration}

The TKG-Transformer and the patient-graph GNN each output a genuine cumulative-incidence probability per disease and horizon, so we assess their calibration directly with a reliability diagram: predicted probabilities are grouped into ten population deciles, and the observed event rate within each decile is compared to the predicted probability range it covers. We also report the Brier score for each disease. Cox and XGBoost, as implemented in this study, output a relative risk score from a proportional-hazards or partial-likelihood objective, not a calibrated probability with a fixed scale; a Brier score is therefore not meaningful for their output as produced here. For these two models, we instead check whether risk is ordered correctly on average: patients are grouped into ten population deciles by predicted risk score, and we check whether the observed event rate increases monotonically from the lowest to the highest decile.

\subsubsection{Fairness and Subgroup Analysis}
\label{sec:fairness}

We compute AUROC at the 3-year horizon separately within subgroups defined by sex, age tertile, and self-reported race/ethnicity category, for all four models. Age tertile cutpoints are computed using training-set patients only and then applied to all splits. We do not apply any predefined minimum sample or event-count threshold to decide which subgroups to report; every subgroup is reported together with its total size and its number of positive cases, so that the reliability of a given subgroup estimate can be judged directly from those counts rather than assumed.
